\documentclass[11pt]{article}

\usepackage[a4paper,margin=1in]{geometry}
\usepackage{amsmath,amssymb}
\usepackage[hidelinks]{hyperref}
\usepackage{enumitem}

\emergencystretch=4em
\sloppy

\title{Fourth Codex Review: Remaining Minor Findings for\\
\texttt{rational\_cut\_and\_project\_gap\_periods.tex}}
\author{Codex Review Notes}
\date{May 13, 2026}

\begin{document}

\maketitle

\section*{Executive Summary}

I reviewed the current final version of
\[
\texttt{LaTeX/rational\_cut\_and\_project\_gap\_periods.tex}
\]
after the previous findings had been implemented.  I do not see a remaining
mathematical blocker.  The rational multiset theorem, the set-valued theorem,
the generic minimality argument, the degenerate case, and the irrational
aperiodicity proof are coherent in the current manuscript.

The previous F3/R2 concern is resolved in the strongest useful way: the paper
now lifts a projected period only to the forward inclusion
\[
  A+v\subseteq A,
\]
and Step~3 derives the contradiction from this forward inclusion alone, in
agreement with the Lean proof.

The points below are therefore only low-importance precision or presentation
items.  They are not blockers for arXiv submission.

\section*{Remaining Points}

\begin{enumerate}[label=\textbf{R\arabic*.},leftmargin=*]

\item \textbf{Replace ``quantor structure'' by ``quantifier structure''.}

Around
\[
\texttt{LaTeX/rational\_cut\_and\_project\_gap\_periods.tex:1267},
\]
the text says:
\begin{quote}
``construction whose quantor structure is the natural one''
\end{quote}
The word ``quantor'' is a Germanism.  In English mathematical prose, the
idiomatic term is ``quantifier''.

\textbf{Recommended wording.}
\begin{quote}
``construction whose quantifier structure is the natural one''
\end{quote}

\textbf{Importance: very low.}  This is purely linguistic.

\item \textbf{Slightly sharpen the local-finiteness wording.}

Around
\[
\texttt{LaTeX/rational\_cut\_and\_project\_gap\_periods.tex:1261},
\]
the phrase
\begin{quote}
``\(\tilde p\) restricted to the strip has finite fibres over any bounded
interval''
\end{quote}
is understandable, but a little nonstandard.  A fibre is usually over a
single value, not over an interval.

\textbf{Recommended wording.}
\begin{quote}
``the preimage of every bounded \(\tilde p\)-interval inside the strip is
finite''
\end{quote}
or equivalently
\begin{quote}
``the intersection of the strip with the preimage of a bounded
\(\tilde p\)-interval is a bounded parallelogram, hence contains only
finitely many lattice points.''
\end{quote}

\textbf{Importance: low.}  The argument is correct as written; this would
only remove a small terminology imprecision.

\item \textbf{Optionally reorder the \(\varepsilon\)-choice paragraph.}

The unboundedness argument around
\[
\texttt{LaTeX/rational\_cut\_and\_project\_gap\_periods.tex:1272--1308}
\]
is mathematically correct.  The only possible reader stumble is that the text
first discusses a density witness \((p_0,q_0)\) for an arbitrary
\(\varepsilon\), then later chooses \(\varepsilon\) small depending on the
target \(M\).

\textbf{Recommended logical order.}
\begin{enumerate}[label=(\roman*),leftmargin=*]
  \item Fix \(M\).
  \item Choose
  \[
    0<\varepsilon\le
    \min\!\left(1,\,
      \frac{c_0\min(a\omega,\omega)}{|M|+2c_0+1}\right).
  \]
  \item Use density to choose \((p_0,q_0)\) with
  \(0<|s(p_0,q_0)|<\varepsilon\).
  \item Define
  \[
    K=\max\!\bigl(1,\lceil M/\tilde p(p_0,q_0)\rceil+1\bigr).
  \]
\end{enumerate}

\textbf{Importance: low.}  This is not a flaw in the proof; it is only a
possible readability improvement for a referee checking the quantifiers.

\item \textbf{Use either ``formalization'' or ``formalisation'' consistently.}

Around
\[
\texttt{LaTeX/rational\_cut\_and\_project\_gap\_periods.tex:1318},
\]
the text uses ``formalisation'', while most of the paper uses the American
spelling ``formalization''.

\textbf{Importance: very low.}  This is only copy-editing consistency.

\item \textbf{Optional typography: two harmless underfull box warnings.}

A fresh \texttt{pdflatex} build reports no errors, no undefined references,
no overfull boxes, and only two underfull box warnings in the Lean/table
area:
\[
\texttt{lines 1456--1473}
\qquad\text{and}\qquad
\texttt{line 1579}.
\]

\textbf{Importance: very low.}  Underfull boxes are usually harmless.  They
need not be fixed unless the final PDF visually looks awkward at those spots.

\end{enumerate}

\section*{Checks Performed}

\begin{itemize}[leftmargin=*]
  \item \texttt{pdflatex} on the current paper completed successfully.
        The searched log showed no undefined references, no citation
        warnings, no errors, and no overfull boxes.
  \item \texttt{lake build} in \texttt{Lean/CutAndProject} completed
        successfully:
        \[
          \texttt{Build completed successfully (8251 jobs).}
        \]
  \item A search of the Lean source found no occurrences of
        \texttt{sorry}, \texttt{admit}, or \texttt{axiom}.
  \item Running
        \[
          \texttt{python3 src/python/verify\_paper.py}
        \]
        reported zero multiset mismatches, zero set mismatches, and zero
        negative set sentinels over all \(51{,}729\) CSV rows.
\end{itemize}

\section*{Overall Assessment}

The manuscript looks arXiv-ready from the mathematical side.  The remaining
items above are polish: they improve English idiom, terminology, and
quantifier readability, but they do not affect the correctness of the main
results.

\end{document}
